\documentclass[a4paper]{article}
\usepackage[pdftex]{graphicx}
\usepackage[utf8]{inputenc}
\usepackage{enumerate}
\usepackage{siunitx}
\sisetup{locale=DE}
\usepackage{href-ul}
\hypersetup{
	colorlinks=true,
	linkcolor=blue,
	urlcolor=blue}
\usepackage{icomma}
\usepackage{amssymb}
\usepackage{geometry}
\geometry{a4paper, top=15mm, left=15mm, right=15mm, bottom=15mm,
	headsep=10mm, footskip=12mm}

\begin{document}
	%Title
	{\bf Electric field strength and Coulomb's law }\\
	\begin{enumerate}[1.]
		
		%Task 1
		\item {\bf Hydrogen atom}
		\begin{enumerate}[a)]
			\item Calculate the magnitude of the electric force $F_C$ between the nucleus and the electron of a hydrogen atom. The following applies to the distance $r$ of the electron from the proton: $r=\SI{5,3e-11}{\meter}$ (=radius of the Bohr basic orbit).
			\item Compare this value $F_C$ with the gravitational force $F_G$ between the nucleus and the electron.
		\end{enumerate}
		
		%Exercise 2
		\item {\bf Cancellation of gravitational force and Coulomb force}\\
		Two bodies of mass $m_1$ and $m_2$ that can be viewed as mass points carry the positive charges $Q_1$ and $Q_2$ with $m_1=m_2=\SI{1,0e-3}{\kilogram} $ and \\
		$Q_1= \SI{1.0e-13}{\coulomb}$. How large does the charge $Q_2$ have to be so that the gravitational force and the Coulomb force cancel each other out?
		
		%Task 3
		\item {\bf Pendulum with two balls charged with the same name}\\
		Two identical balls with the weight of $F_G=\SI{0.5e-2}{\newton}$ are attached to two each
		$l=\SI{1.00}{\meter}$ long threads attached to the same point at the top and carry equal charges $Q_1=Q_2=Q$. Because of the repulsion, the balls have a distance of $d=\SI{0.40}{\meter}$. Calculate the amount of charge on the balls.
		
		%Task 4
		\item {\bf Positively charged metal ball}\\
		A metal ball of radius $R=\SI{4.5}{\centi\meter}$ carries a positive charge \\
		$Q=\SI{1.0e-9}{\coulomb}$.
		\begin{enumerate}[a)]
			\item Specify the magnitude of the electric field strength with inserted numerical values depending on the distance $r$.
			\item Calculate the magnitude of the electric field strength in the range \\ $R \le r \le 3 R$.
			\item Graphically represent the course of the electric field strength for $R \le r \le 3 R$.
			\item Specify the magnitude of the electric field strength inside the sphere and complete the diagram in subtask c) accordingly.
			\item Calculate the force on a helium nucleus ($\alpha$ particle) located at a distance $r_{\alpha}=\SI{13.5}{\centi\meter}$ from the center of the sphere.
		\end{enumerate}
		
		%Task 5
		\item {\bf balloon}\\
		A spherical balloon is coated with conductive graphite on the outside and suspended in isolation. Its radius $r_1$ is $\SI{2.5}{\centi\meter}$. The charge $Q$ is evenly distributed on the balloon. At a distance $d=\SI{52.5}{\centi\meter}$ from the center of the sphere, this charge creates the electric field strength $E_1=\SI{100}{\newton\over \coulomb}$.
		\begin{enumerate}[a)]
			\item Calculate the amount of charge $Q$.
			\item The balloon is now inflated, the new radius $r_2$ is $\SI{5,0}{\centi\meter}$. Determine the field strength $E_2$, which occurs at the distance $d$ from the center of the sphere.
		\end{enumerate}
		
		%Task 6
		\item {\bf Field strength-free point}\\
		Two positive charges are given with $Q_1=\SI{8.5}{\nano\coulomb}$ and $Q_2=\SI{5.5}{\nano\coulomb}$. On the line connecting the two charges there is a point $P$ at a distance $s=\SI{10,0}{\centi\meter}$ from the charge $Q_1$ in which the electric field strength is zero. Calculate the distance $d$ between the two point charges.
	\end{enumerate}
	
	\fbox{
		\begin{minipage}{0.5\textwidth}
			For the solution, please \href{https://www.okuyakl.de/phys/p11coulL240/le240.pdf}{click here} or scan the QR code.\\
			Additional worksheets are available at
			
			\href{http://www.okuyakl.com}{www.okuyakl.com}
		\end{minipage}
		\hfill
		\begin{minipage}{0.4\textwidth}
			\includegraphics[width=1.5 cm]{../../viecher/zwe03}
			\includegraphics[width=3 cm]{qre240}
			\includegraphics[width=2 cm]{../../viecher/afanticon1}
			
	\end{minipage}}
	
\end{document}